\exercisegroup

\begin{enumerate}
\def\labelenumi{\arabic{enumi})}
\tightlist
\item
  设 A 是 \(\mathbf{R}^2\) 的闭离散子集。
\end{enumerate}

\begin{enumerate}
\def\labelenumi{\alph{enumi})}
\item
  对于 R \textgreater{} 0，令 \(D_{R}\) 为这样的 \(z \in C\) 所成的集合，其中 \(|z| < R\)；证明 \(D_{R} - (A \cap D_{R})\) 的基本群同构于以 \(A \cap D_{R}\) 为基的自由群。
\item
  证明存在同胚 \(f\)，它把 \(\mathbf{R}^{2}\) 映入自身，并满足 \(f(\mathrm{A})\subset\mathbf{N}\times\{0\}\) 。
\item
  证明 \(\mathbf{R}^2 - \mathbf{A}\) 的基本群同构于 \(\mathrm{F(A)}\) 。
\end{enumerate}

\begin{enumerate}
\def\labelenumi{\arabic{enumi})}
\setcounter{enumi}{1}
\item
  设 A 是 \(R^{2}\) 中坐标为 \((1/n,0)\) 的点集，对 \(n\in N^{*}\) 而言。证明 \(R^{2}-A\) 的基本群具有连续统的基数。特别地，它不同构于群 F(A)。
\item
  设 \(((\mathrm{X}_{i}, x_{i}))_{i \in \mathrm{I}}\) 是带基点拓扑空间族；记 \((\mathrm{X}, x)\) 为该族的楔和（IV，第 421 页，例 4）。假设对每个 \(i \in I\)，点 \(x_{i}\) 在 \(X_{i}\) 中是闭点，并且具有邻域基 V，使得群 \(\pi_{1}(\mathrm{V}, x_{i})\) 归结为单位元。则典范同态 \(*_{i \in \mathrm{I}} \pi_{1}(\mathrm{X}_{i}, x_{i}) \to \pi_{1}(\mathrm{X}, x)\) 是同构。
\item
  设 P 是如 I，第 146 页练习 5 b) 中所定义的 \(R^{2}\) 子空间。令 a 为点 \((0,0,1)\)（它属于 \(R^{3}\)），令 C 为线段 \([a,(x,0)]\)（其中 \(x \in P\)）的并。
\end{enumerate}

\begin{enumerate}
\def\labelenumi{\alph{enumi})}
\item
  证明空间 C 局部道路连通且道路单连通。
\item
  设 D 是空间 C 与其对称像 -C 的并。证明空间 D 单连通，但不道路单连通。
\item
  证明庞加莱群 \(\pi_{1}(D,0)\) 的基数为连续统的基数。
\end{enumerate}

\begin{enumerate}
\def\labelenumi{\arabic{enumi})}
\setcounter{enumi}{4}
\tightlist
\item
  设 P 是 \(R^{2}\) 的子空间，如 I 第 146 页练习 5 b) 所定义；另记 \(P^{+}\)（分别为 \(P^{-}\)）为点 \((x,y)\in\mathrm{P}\) 的集合，其中 \(y\geqslant0\)（分别为 \(y\leqslant0\)）。设 a 是点 \((0,0,1)\)（位于 \(R^{3}\) 中），设 A 是线段 [a,0] 与各线段 \([a,(2r_{n},0)]\)（\(n\in N\)）的并。置 \(B^{+}=(P^{+}\times\{0\})\cup A\)，\(B^{-}=(P^{-}\times\{0\})\cup A\)，且 \(B=(P\times\{0\})\cup A\)。
\end{enumerate}

\begin{enumerate}
\def\labelenumi{\alph{enumi})}
\item
  证明空间 \(B^{+}\) 可解圈，且庞加莱群 \(\pi_{1}(B^{+}, a)\) 是以 N 为基的自由群。
\item
  证明 A 可缩到 a。
\item
  证明从 \(\pi_{1}(B^{+},a)*\pi_{1}(B^{-},a)\) 到 \(\pi_{1}(B,a)\) 的典范同态不是满射；它由 \(B^{+}\) 和 \(B^{-}\) 到 B 的典范包含诱导。
\end{enumerate}

\begin{enumerate}
\def\labelenumi{\arabic{enumi})}
\setcounter{enumi}{5}
\tightlist
\item
  设 X 是实直线 R，A 是 X 的子集。设 Y 是将 A 缩并为一点后由 X 得到的空间。
\end{enumerate}

\begin{enumerate}
\def\labelenumi{\alph{enumi})}
\item
  若 A 离散，则 Y 的庞加莱群是自由群。
\item
  若 A 有聚点，则 Y 的庞加莱群的基数为连续统的基数，且不是自由群。
\item
  若 A 是点集 \(1/n\)（\(n \in \mathbf{N}^*\)）的闭包，则空间 X/A 同胚于夏威夷耳环。
\end{enumerate}

\begin{enumerate}
\def\labelenumi{\arabic{enumi})}
\setcounter{enumi}{6}
\tightlist
\item
  设 G 是图，S 是其顶点集，A 是其定向边集。记 \(|G|\) 为其几何实现；将 S 识为 \(|G|\) 的子集，并记 \(p: I \times A \rightarrow |G|\) 为典范映射。
\end{enumerate}

\begin{enumerate}
\def\labelenumi{\alph{enumi})}
\item
  证明空间 \(|G|\) 紧（分别局部紧）当且仅当图 G 有限（分别局部有限）。
\item
  证明从 S 到 \(|\mathbf{G}|\) 的典范映射在 \(\pi_0(\mathbf{G})\) 与 \(\pi_0(|\mathbf{G}|)\) 之间诱导双射。
\item
  记 \(\varpi(|\mathrm{G}|)_\mathrm{S}\) 为 \(\varpi(|\mathrm{G}|)\) 的全子群胚，其顶点集为 S。证明存在唯一同构 \(\varpi (\mathrm{G})\to \varpi (|\mathrm{G}|)_{\mathrm{S}}\)，它在顶点集上为恒等，并将定向边 \(a\in \mathrm{A}\) 的道路类映为道路 \(t\mapsto p(t,a)\) 的类
\item
  设图 G 连通且有限；令 \(m = 1 + \text{Card}(S) - \text{Card}(A)\)。证明 \(|G|\) 的庞加莱群是有 m 个生成元的自由群。
\end{enumerate}

\begin{enumerate}
\def\labelenumi{\arabic{enumi})}
\setcounter{enumi}{7}
\tightlist
\item
  设 X 是空间 \(R \times \{0,1\}\) 的商空间，商所取的是使 \((x,0)\) 与 \((x,1)\) 对所有 \(x \in R^{*}\) 均等价的最粗等价关系（“原点加倍的实直线”）。
\end{enumerate}

\begin{enumerate}
\def\labelenumi{\alph{enumi})}
\item
  证明 \(\pi_1(\mathrm{X},x)\) 同构于 \(\mathbf{Z}\)，对每个点 \(x\in \mathrm{X}\) 均如此。
\item
  设 \(\mathrm{X}^{(2)}\) 是空间 \(X^{2}\) 在对称群 \(S_{2}\) 的作用下所得的商空间，其中该作用通过置换坐标于空间 \(X^{2}\) 上进行。证明 \(\pi_{1}(\mathrm{X}^{(2)},y)=0\)，对所有 \(y\in\mathrm{X}^{(2)}\) 均成立。
\end{enumerate}

\begin{enumerate}
\def\labelenumi{\arabic{enumi})}
\setcounter{enumi}{8}
\tightlist
\item
  设 X 是拓扑空间，设 R 是 X 上的等价关系，设 \(p: X \to X/R\) 是典范满射。设 \(x \in X\)。
\end{enumerate}

\begin{enumerate}
\def\labelenumi{\alph{enumi})}
\item
  设 X 连通且局部道路连通，等价类是连通集，且空间 X/R 可解圈。证明同态 \(p_{*}:\pi_{1}(X,x)\to\pi_{1}(X/R,p(x))\) 是满射。
\item
  设 X = I，且 \(\{0,1\}\) 是唯一不退化为单点的等价类。证明同态 \(p_{*}\) 不是满射。
\end{enumerate}

\begin{enumerate}
\def\labelenumi{\arabic{enumi})}
\setcounter{enumi}{9}
\tightlist
\item
  设 W 是华沙圆，定义见 IV 第 455 页练习 1，现重新采用其记号。
\end{enumerate}

\begin{enumerate}
\def\labelenumi{\alph{enumi})}
\item
  设 R 是 W 上的等价关系，其唯一的非单点等价类为 B 和 S；设 \(p: W \rightarrow W/R\) 为典范满射。证明商空间 W/R 同伦等价于圆但不同胚于圆 \(S_{1}\)。证明同态 \(p_{*}\) 不是满射。
\item
 设 R' 是 W 上的等价关系，使 B ∪ S 成为唯一不退化为单点的等价类；设 \(p'\) : W → W/R' 为典范满射。证明空间 W/R' 同胚于 \(S_{1}\)，同态 \(p_{*}\) 不是满射，但 R' 的等价类是连通的。
\end{enumerate}

\begin{enumerate}
\def\labelenumi{\arabic{enumi})}
\setcounter{enumi}{10}
\tightlist
\item
  设 X 是道路连通拓扑空间，设 \(f: X \to X\) 是同胚。设 T 是 \(X \times I\) 的商空间，商所取的是使点 \((x,0)\) 与 \((f(x),1)\) 对每个 \(x \in X\) 均等价的最细等价关系。记 \(p: X \times I \to T\) 为典范映射；记 \(j\colon X\to T\) 为由 \(x\mapsto p(x,1/2)\) 给出的映射。设 a 是 X 中的一点，设 \(c\in\Lambda_{f(a),a}(X)\) 是起点为 \(f(a)\)、终点为 \(a\) 的道路，并记 \(\gamma=[c]\)。
\end{enumerate}

\begin{enumerate}
\def\labelenumi{\alph{enumi})}
\item
  证明从 I 到 T 的映射由 \(t \mapsto p(a, \frac{1}{2} - t)\)（当 \(t \in [0, \frac{1}{2}]\)）以及 \(t \mapsto p(c(2t - 1), \frac{3}{2} - t)\)（当 \(t \in ]\frac{1}{2}, 1]\)）给出，是 T 中基于 \(j(a)\) 的圈。记其类为 \(\delta\)。
\item
 设 S 是从 \(\pi_{1}(\mathrm{X},a)*\mathbf{Z}\) 到 \(\pi_{1}(\mathrm{T},j(a))\) 的唯一群同态，它将类 \(v\in\pi_{1}(\mathrm{X},a)\) 映为 \(j_{*}(v)\)，并将 Z 中的元素 t=1 映为类 \(\delta\)。证明同态 S 是满射，且其核是包含元素 \(vt(\gamma^{-1}f_{*}(v)\gamma)^{-1}t^{-1}\)（\(v\in\pi_{1}(\mathrm{X},a)\)）的最小正规子群。
\item
  对 \(v \in \pi_{1}(X, a)\)，置 \(\varphi(v) = \gamma^{-1} f_{*}(v) \gamma\)；证明 \(\varphi\) 是群 \(\pi_{1}(X, a)\) 的自同构。设 \(\varphi: \mathbf{Z} \to \operatorname{Aut}(\pi_{1}(X, a))\) 是将 1 映为 \(\varphi\) 的唯一群同态。构造从群 \(\pi_{1}(T, j(a))\) 到 Z 与 \(\pi_{1}(X, a)\) 的外半直积（对应作用为 \(\varphi\)）的同构。
\end{enumerate}

\begin{enumerate}
\def\labelenumi{\arabic{enumi})}
\setcounter{enumi}{11}
\tightlist
\item
  \begin{enumerate}
  \def\labelenumii{\alph{enumii})}
  \tightlist
  \item
  设 D 是 \(R^{3}\) 中的一条直线，令 \(U = R^{3} - D\)。证明 U 道路连通，并且对每个点 \(a \in U\)，群 \(\pi_{1}(U, a)\) 同构于 Z。
  \end{enumerate}
\end{enumerate}

\begin{enumerate}
\def\labelenumi{\alph{enumi})}
\setcounter{enumi}{1}
\tightlist
\item
  设 P 是 \(R^{3}\) 中的一个平面，C 是包含于 P 的圆，令 \(V = R^{3} - C\)。证明 V 道路连通，且对每个点 \(a \in V\)，群 \(\pi_{1}(V, a)\) 同构于 Z。（将范坎彭定理应用于由平面 P 界定的两个半空间构成的 \(R^{3}\) 的闭覆盖；也可以注意到空间 U 和 V 同胚。）
\end{enumerate}

\begin{enumerate}
\def\labelenumi{\arabic{enumi})}
\setcounter{enumi}{12}
\tightlist
\item
    设 D 和 D' 是 \(R^{3}\) 中两条不同的直线，令 \(U = R^{3} - (D \cup D')\)。
\end{enumerate}

\begin{enumerate}
\def\labelenumi{\alph{enumi})}
\item
  证明 U 是道路连通的。
\item
  假设 D 和 D' 不相交；证明对每个点 \(a \in U\)，群 \(\pi_{1}(U, a)\) 是以两个生成元为基的自由群。
\item
  假设 D 和 D' 有公共点；证明对每个点 \(a \in U\)，群 \(\pi_{1}(U, a)\) 是以三个生成元为基的自由群。
\end{enumerate}

\begin{enumerate}
\def\labelenumi{\arabic{enumi})}
\setcounter{enumi}{13}
\tightlist
\item
    设 P 是 \(R^{3}\) 中的一个平面，C 是包含于 P 的圆，D 是 \(R^{3}\) 中的一条直线。令 \(U = R^{3} - (C \cup D)\)，并令 a 为 U 中的一点。
\end{enumerate}

\begin{enumerate}
\def\labelenumi{\alph{enumi})}
\item
  证明 U 是道路连通的。
\item
  假设 D 不与 C 的凸包相交。证明 \(\pi_{1}(U,a)\) 是以两个生成元为基的自由群。
\item
  假设 D 不与 C 相交，但与 C 的凸包相交。证明 \(\pi_{1}(U,a)\) 同构于 \(Z^{2}\)。
\item
  假设 D 与 C 恰好交于一点。证明 \(\pi_{1}(U,a)\) 是以两个生成元为基的自由群。
\item
  假设 D 与 C 交于两点。证明 \(\pi_{1}(U,a)\) 是以三个生成元为基的自由群。
\end{enumerate}

\begin{enumerate}
\def\labelenumi{\arabic{enumi})}
\setcounter{enumi}{14}
\tightlist
\item
  若拓扑空间 X 是连通的，并且对于 X 的每一对不交闭子集 (A, B)，只要 \(X-A\) 和 \(X-B\) 连通，就有 \(X-(A \cup B)\) 连通，则称 X 具有弗拉格门–布劳威尔性质。
\end{enumerate}

\begin{enumerate}
\def\labelenumi{\alph{enumi})}
\item
  设 X 是具有弗拉格门–布劳威尔性质的局部连通拓扑空间。令 (A, B) 是 X 的一对不交闭子集，令 \((a, b)\) 是 \(X-(A \cup B)\) 中的一对点，它们属于 \(X-A\)（分别属于 \(X-B\)）的同一连通分支。证明 a、b 属于 \(X-(A \cup B)\) 的同一连通分支。
\item
  设 X 是连通且局部道路连通的拓扑空间，但不具有弗拉格门–布劳威尔性质。证明，对每个点 \(x \in X\)，存在从 \(\pi_{1}(X, x)\) 到 Z 的满射同态。
\item
  设 X 是连通、局部道路连通且道路单连通的拓扑空间。证明 X 具有弗拉格门–布劳威尔性质。
\end{enumerate}

d) 设 X 是分离且局部连通的拓扑空间，满足对每个点 \(x \in \mathrm{X}\)，空间 \(\mathrm{X} - \{x\}\) 具有弗拉格门–布劳威尔性质。令 A 是同胚于 \([0,1]\) 的 X 的子集；证明 \(\mathrm{X} - \mathrm{A}\) 连通。（令 \(c: \mathbf{I} \to \mathrm{A}\) 为同胚。反设并构造一个序列 \((\mathrm{J}_n)\)，使得对每个 \(n\)，\(\mathrm{J}_n\) 是 \(\mathbf{I}\) 中长度为 \(2^{-n}\) 的区间，包含于 \(\mathrm{J}_{n-1}\) 中（若 \(n \geqslant 1\)），并且 \(\mathrm{X} - c(\mathrm{J}_n)\) 不连通；令 \(x\) 为 X 中属于所有集合 \(c(\mathrm{J}_n)\) 的唯一点。然后证明 \(\mathrm{X} - \{x\}\) 不连通。）

\begin{enumerate}
\def\labelenumi{\alph{enumi})}
\setcounter{enumi}{4}
\tightlist
\item
  设 \(n\) 是整数且 \(\geqslant 1\)，令 A 是 \(\mathbf{S}_n\) 中同胚于 \([0,1]\) 的子集。证明 \(\mathbf{S}_n - \mathbf{A}\) 连通。
\end{enumerate}

\begin{enumerate}
\def\labelenumi{\arabic{enumi})}
\setcounter{enumi}{15}
\tightlist
\item
  设 A 是 \(S_{2}\) 中同胚于 \(S_{1}\) 的子集。令 a、b 是 A 的不同点；记 \(A'\) 和 \(A''\) 为 \(A-\{a,b\}\) 的连通分支，并置 \(X=S_{2}-\{a,b\}\)、\(U=X-A'\) 和 \(V=X-A''\)。
\end{enumerate}

\begin{enumerate}
\def\labelenumi{\alph{enumi})}
\item
  证明 U 和 V 连通。（使用练习 15。）
\item
  证明对每个点 \(x \in U\)，由 U 包含于 X 所诱导的从 \(\pi_{1}(U, x)\) 到 \(\pi_{1}(X, x)\) 的同态是平凡的。
\item
  证明 \(S_{2}-A\) 恰有两个连通分支，且它们的边界等于 A。
\item
  设 S 是 \(R^{2}\) 中同胚于 \(S_{1}\) 的子集；证明 \(R^{2}-S\) 恰有两个连通分支，且它们的边界等于 S（若尔当定理）。
\end{enumerate}

\begin{enumerate}
\def\labelenumi{\arabic{enumi})}
\setcounter{enumi}{16}
\tightlist
\item
  我们将 \(S_{1}\) 识别为模长为 1 的复数集，并将球面 \(S_{3}\) 识别为 \(C^{2}\) 的子空间，由点 \((z,w)\) 组成，这些点满足 \(|z|^{2} + |w|^{2} = 1\)。令 B、C 分别为 \(S_{3}\) 中由 \(|z| \leqslant |w|\) 和 \(|z| \geqslant |w|\) 定义的子空间。
\end{enumerate}

\begin{enumerate}
\def\labelenumi{\alph{enumi})}
\item
  证明 B 和 C 同胚于 \(B_{2} \times S_{1}\)。
\item
  对于每一对互素整数 \((m,n)\)，记 \(K_{m,n}\) 为从 \(S_{1}\) 到 \(S_{3}\) 的映射 \(u\mapsto(u^{m},u^{n})\) 的像（“环面纽结”）。
\item
  证明 \(S_{3}-K_{m,n}\) 连通，且其基本群具有表现 \(\langle x,y;x^{m}=y^{n}\rangle\)。
\end{enumerate}

d) 证明，存在同胚 \(\varphi\)，它从 \(\mathbf{S}_3\) 到自身，并且满足 \(\varphi(\mathrm{K}_{m,n}) = \mathrm{K}_{1,0}\)，当且仅当 \(|m| \leqslant 1\) 或 \(|n| \leqslant 1\)。

\begin{enumerate}
\def\labelenumi{\arabic{enumi})}
\setcounter{enumi}{17}
\tightlist
\item
  设 \((m,n)\) 是满足 1 \textless{} m \textless{} n 的一对自然整数。令 G 为由表现 \(\langle x,y;x^{m}=y^{n}\rangle\) 定义的群。令 \(z=x^{m}\)，令 C 为 G 中由 z 生成的子群。
\end{enumerate}

\begin{enumerate}
\def\labelenumi{\alph{enumi})}
\item
  证明 C 包含于 G 的中心。
\item
  证明 G/C 同构于自由积 \((\mathbf{Z}/m\mathbf{Z})*(\mathbf{Z}/n\mathbf{Z})\)。
\item
  推出 C 是 G 的中心。
\item
  设 \((p,q)\) 是满足 1 \textless{} p \textless{} q 的一对自然整数。假设群 \(G'\) 由表现 \(\langle x,y; x^{p} = y^{q}\rangle\) 定义且同构于 G。证明 \((p,q) = (m,n)\)。
\item
  证明 G 对其导出群的商同构于 Z。
\end{enumerate}

\begin{enumerate}
\def\labelenumi{\arabic{enumi})}
\setcounter{enumi}{18}
\tightlist
\item
  设 G 是图；记 S 为其顶点集，A 为其有向边集，\(|\mathbf{G}|\) 为其几何实现。记 \(j\colon \mathrm{S}\to |\mathrm{G}|\) 和 \(p\colon \mathbf{I}\times \mathbf{A}\to |\mathbf{G}|\) 为典范映射。假设集合 S 和 A 都是有限的。
\end{enumerate}

一个映射 \(f\colon|G|\to\mathbf{R}^{n}\) 称为仿射的（分别为分段仿射的），若对每条边 \(a\in A\)，映射 \(t\mapsto f\circ p(a,t)\) 是仿射的（分别为分段仿射的）。（参见 III，第 331 页，练习 6。）

\begin{enumerate}
\def\labelenumi{\alph{enumi})}
\item
   证明，存在从 \(|G|\) 到 \(R^{3}\) 的单射分段仿射映射。（先证明存在整数 \(n \geqslant 1\) 以及从 \(|G|\) 到 \(R^{n}\) 的单射分段仿射映射。）
\item
   对每个连续单射映射 \(f\colon|G|\to R^{n}\)，证明存在连续单射映射 \(g\colon|G|\to R^{n}\)，它是分段仿射的并且与 f 同伦。
\item
   假设映射 \((o,t)\colon\mathrm{A}\to\mathrm{S}\times\mathrm{S}\) 是单射。令 f 为从 S 到 R 的单射映射。证明，存在唯一的从 \(|G|\) 到 \(R^{3}\) 的仿射映射 F，它把每个顶点 \(s\in S\) 映到点 \((f(s),f(s)^{2},f(s)^{3})\)（该点属于 \(R^{3}\)）。证明 F 是单射。
\end{enumerate}

我们称图 G 是平面的，如果存在从 \(|G|\) 到 \(R^{2}\) 的连续单射映射。

d) 假设图 \(\mathbf{G}\) 是平面的，且映射 \((o,t)\colon\mathbf{A}\to\mathbf{S}\times\mathbf{S}\) 是单射。令 \(f\colon|\mathbf{G}|\to\mathbf{R}^{2}\) 是连续单射映射。证明，存在单射仿射映射 \(g\colon|\mathbf{G}|\to\mathbf{R}^{2}\)，它与 \(f\) 同伦。

\begin{enumerate}
\def\labelenumi{\arabic{enumi})}
\setcounter{enumi}{19}
\tightlist
\item
   设 \(f \colon \mathbf{I} \to \mathbf{R}^2\) 是一个分段仿射圈，且其在 [0,1[ 上的限制是单射；记 \(C = f(\mathbf{I})\)。
\end{enumerate}

\begin{enumerate}
\def\labelenumi{\alph{enumi})}
\item
   证明 \(\mathbf{R}^2 -\mathbf{C}\) 至多有两个连通分支。
\item
   对于 \(x \in R^{2} - C\) 和 \(v \in S_{1}\)，令 \(n_{v}(x)\) 为 \((\mathbf{R}^{2} - \mathbf{C}) \cap (x + \mathbf{R}_{+}v)\) 的连通分支数。证明，存在唯一的局部常值映射 \(n: R^{2} - C \to \{0,1\}\)，使得 \(n(x) \equiv n_{v}(x) \pmod{2}\) 对所有 \(v \in S_{1}\) 都成立。
\item
   推出 \(R^{2}-C\) 恰有两个连通分支，且它们的边界等于 C。
\item
   设 \(g: I \to R^{2}\) 是单射分段仿射映射，满足 \(\overline{g}^{-1}(C) = \{0,1\}\)；置 \(P = g(I)\)，\(D = C \cup P\)。证明 \(C - (C \cap P)\) 有两个连通分支；记它们为 \(P_{1}\) 和 \(P_{2}\)。证明 \(R^{2} - D\) 有三个连通分支，且它们的边界分别等于 C、\(P_{1} \cup P\) 和 \(P_{2} \cup P\)。
\item
   设 G 是有限平面图；令 S 为其顶点集，A 为其定向边集。令 \(f\colon|G|\to\mathbf{R}^{2}\) 是连续单射分段仿射映射，令 P 为其像。证明 \(\mathbf{R}^{2}-P\) 的连通分支数等于 \(1+\mathrm{Card(S)}-\mathrm{Card(A)}+\mathrm{Card}(\pi_{0}(\mathrm{G}))\)。
\end{enumerate}

\begin{enumerate}
\def\labelenumi{\arabic{enumi})}
\setcounter{enumi}{20}
\item
  令 K 为与如下箭图相关的图：其顶点集为 \(\{1,2,3\} \times \{0,1\}\)，箭集为形如 \(((a,0),(b,1))\) 的所有对，其中 \(a,b\in \{1,2,3\}\)，而其源映射和靶映射分别由第一、第二投影导出（“库拉托夫斯基图”）。证明图 K 不是平面的。
\item
  设 \(n\) 是自然数，令 \(\mathrm{K}_n\) 为顶点集基数为 \(n\) 的完全图（II，第 220 页，练习 5）。
\end{enumerate}

\begin{enumerate}
\def\labelenumi{\alph{enumi})}
\item
  证明图 \(K_{n}\) 在 \(1 \leqslant n \leqslant 4\) 时是平面的。
\item
  证明图 \(K_{5}\) 不是平面的。
\end{enumerate}

\begin{enumerate}
\def\labelenumi{\arabic{enumi})}
\setcounter{enumi}{22}
\tightlist
\item
   设 \(u\)、\(v\) 是 \(\mathbf{S}_1\) 中满足 \(\mathcal{I}(v / u) > 0\) 的两个点；定义族 \((a_j)_{j \in \mathbf{Z} / 4\mathbf{Z}}\)，其中 \(a_1 = u\)、\(a_2 = v\)、\(a_3 = -u\)、\(a_4 = -v\)。记 \(\mathrm{Q}_1\)、\(\mathrm{Q}_2\) 为线段 \([a_1, a_3]\) 和 \([a_2, a_4]\)，它们位于单位圆盘 \(\mathbf{B}_2\) 中，并置 \(\mathrm{Q} = \mathrm{Q}_1 \cup \mathrm{Q}_2\)。
\end{enumerate}

\begin{enumerate}
\def\labelenumi{\alph{enumi})}
\tightlist
\item
  证明，对每个 \(j \in Z/4Z\)，存在唯一的 \(A_{j} \in \pi_{0}(B_{2}-Q)\) 的道路连通分支，其闭包包含 \(a_{j}\) 和 \(a_{j-1}\)。证明族 \((A_{j})\) 的并等于 \(B_{2}-Q\)。
\end{enumerate}

  置 \(C = B_{2} \times [-1, 1]\)、\(N = (0, 0, 1)\)、\(S = (0, 0, -1)\)。令 \(f: [-1, 1] \to R_{+}\) 为连续映射，满足 \(f(-1) = f(1) = 0\)，且 \(0 < f(t) < 1\) 对所有 \(t \in ]-1; 1[\) 都成立。于是记 \(L_{1}\)、\(L_{2}\) 为从 \([-1, 1]\) 到 C 的映射 \(t \mapsto (ta_{1}, 0)\) 和 \(t \mapsto (ta_{2}, f(t))\) 的像；最后置 \(L = L_{1} \cup L_{2}\)。

\begin{enumerate}
\def\labelenumi{\alph{enumi})}
\setcounter{enumi}{1}
\item
  证明 C-L 道路连通。
\item
  若一条道路 \(c\) 位于 \(\mathbf{C} - \mathbf{L}\) 中，起点为 \(\mathbf{N}\) 且终点为 \(\mathbf{S}\)，并且存在相邻道路 \(c_1\)、\(c_2\) 位于 \(\mathbf{B}_2\) 中，使得 \(c\) 是由路径 \(t \mapsto (c_1(t), 1)\)、\(t \mapsto (c_1(1), 1 - 2t)\) 和 \(t \mapsto (c_2(t), -1)\) 拼接而成的，则称 c 为直的。证明此时 \(c_1(1) \not\in \mathbb{Q}\)；记 \(\mathrm{A}_c\) 为包含该点的 \(\mathbf{B}_2 - \mathbf{Q}\) 的道路连通分支。
\end{enumerate}

d) 设 \(c\) 和 \(c'\) 是 \(\mathrm{C} - \mathrm{L}\) 中起点为 \(\mathrm{N}\)、终点为 \(\mathrm{S}\) 的道路；假设它们是直的。证明它们严格同伦，当且仅当 \(\mathrm{A}_c = \mathrm{A}_{c'}\)。

\begin{enumerate}
\def\labelenumi{\alph{enumi})}
\setcounter{enumi}{4}
\tightlist
\item
  对每个 \(j \in \mathbf{Z}/4\mathbf{Z}\)，令 \(c_j\) 是 \(C-L\) 中起点为 \(N\)、终点为 \(S\) 的直道路，且满足 \(A_{c_j} = A_j\)。证明关系式
\end{enumerate}

\[
[ c _ {1} ] [ c _ {4} ] ^ {- 1} = [ c _ {2} ] [ c _ {3} ] ^ {- 1}
\]

在群 \(\pi_1(\mathrm{C} - \mathrm{L},\mathrm{N})\)

\begin{enumerate}
\def\labelenumi{\alph{enumi})}
\setcounter{enumi}{5}
\tightlist
\item
   证明，从以 \(x, y\) 为两个生成元的自由群到 \(\pi_1(\mathrm{C} - \mathrm{L}, \mathrm{N})\) 的典范同态，将 \(x\) 和 \(y\) 分别映到 \([c_1][c_2]^{-1}\) 和 \([c_2][c_3]^{-1}\)，是群同构。（将范坎彭定理应用于 \(\mathrm{C} - \mathrm{L}\) 的覆盖，该覆盖由 \(\mathbf{R}^3\) 中包含集合 \(\{a_1, a_2, \mathrm{N}\}\) 中两点的三个向量平面所界定的八个闭集组成。）
\end{enumerate}

\begin{enumerate}
\def\labelenumi{\arabic{enumi})}
\setcounter{enumi}{23}
\tightlist
\item
   设 m 是 \(\geqslant 1\) 的整数，令 \((\theta_{j})_{1\leqslant j\leqslant m}\) 是实数序列，满足 \(0\leqslant\theta_{1}<\theta_{2}<\cdots<\theta_{n}<2\pi\)。对 \(j\in Z/mZ\)，置 \(a_{j}=\exp(i\theta_{k})\)，其中 k 是属于类 j 的 \(\{1,\ldots,n\}\) 中唯一元素；再记 \(Q_{j}\) 为线段 \([0,a_{j}]\)，它位于 \(B_{2}\) 中。最后置 \(Q=Q_{1}\cup\cdots\cup Q_{m}\)。
\end{enumerate}

\begin{enumerate}
\def\labelenumi{\alph{enumi})}
\item
  对 \(j \in Z/mZ\)，证明存在唯一的 \(A_{j}\)，它是 \(B_{2}-Q\) 的道路连通分支，并包含所有 \(b \in S_{1}\)，其中满足 \(\mathcal{I}(b/a_{j}) > 0\) 且 \(\mathcal{I}(b/a_{j+1}) < 0\)。证明族 \((A_{j})\) 的成员两两不交，且其并等于 \(B_{2}-Q\)。
\item
   置 \(C = B_{2} \times [-1, 1]\)、\(N = (0, 0, 1)\)、\(S = (0, 0, -1)\)，\(L = Q \times \{0\}\)。若 \(C-L\) 中一条起点为 N、终点为 S 的道路 c 满足存在相邻的道路 \(c_{1}\)、\(c_{2}\) 位于 \(B_{2}\) 中，使得 c 是由路径 \(t \mapsto (c_{1}(t), 1)\)、\(t \mapsto (c_{1}(1), 1 - 2t)\) 和 \(t \mapsto (c_{2}(t), -1)\) 拼接而成的，则称其为直的。证明此时 \(c_{1}(1) \notin Q\)；记 \(A_{c}\) 为包含该点的 \(B_{2}-Q\) 的道路连通分支。
\item
   令 c 和 \(c'\) 是 \(C-L\) 中起点为 N、终点为 S 的道路；假设它们是直的。证明它们严格同伦，当且仅当 \(A_{c} = A_{c'}\)。
\end{enumerate}

d) 对每个 \(j \in \mathbf{Z}/m\mathbf{Z}\)，令 \(c_j\) 是 \(C-L\) 中起点为 N、终点为 S 的直圈，且满足 \(\mathrm{A}_{c_j} = \mathrm{A}_j\)。证明从自由群 \(\mathrm{F}(\mathbf{Z}/m\mathbf{Z})\) 到 \(\pi_1(\mathrm{C}-\mathrm{L},\mathrm{N})\) 的同态将 \(x_j\) 映为 \([c_j][c_{j+1}]^{-1}\)，对每个 \(j \in \mathbf{Z}/m\mathbf{Z}\)，它是满射，且其核是包含乘积 \(x_1x_2 \ldots x_m\) 的最小正规子群。

\begin{enumerate}
\def\labelenumi{\arabic{enumi})}
\setcounter{enumi}{24}
\tightlist
\item
   置 \(C = B_{2} \times [-1, 1]\)。令 K 是 \(R^{3}\) 的闭子空间，令 P 为它在映射 \((x, y, z) \mapsto (x, y)\) 下的像。假设存在有限集 I、点族 \((p_{i})_{i \in \mathrm{I}}\)（点属于 \(R^{2}\)）、严格正实数族 \((r_{i})_{i \in \mathrm{I}}\) 以及自然数族 \((m_{i})_{i \in \mathrm{I}}\)，满足以下性质，其中 \(\varphi_{i}\) 表示从 \(R^{3}\) 到 \(R^{3}\) 的映射 \(x \mapsto (p_{i}, 0) + r_{i}x\)，并且 \(C_{i} = \varphi_{i}(C) \)：
\end{enumerate}

- 各集合 \(C_{i}\) 两两不交；

- 对每个 \(i \in \mathrm{I}\)，集合 \(\mathrm{L}_i = \overline{\varphi}_i^{-1}(\mathrm{C}_i \cap \mathrm{K})\) 是练习 23 中定义的子空间 \(\mathrm{L}\)，若 \(m_i = 0\)，以及练习 24（其中置 \(m = m_i\)），若 \(m_i \geqslant 1\)。

- 集合 \(\mathrm{K}' = \mathrm{K} - \bigcup_{i} (\mathring{\mathrm{C}}_i \cap \mathrm{K})\) 是有限族闭线段的并，且这些线段在平面 \(\mathbf{R}^2 \times \{0\}\) 中两两不交。

令 \(\mathrm{N} = (x_{\mathrm{N}}, y_{\mathrm{N}}, z_{\mathrm{N}})\) 和 \(\mathrm{S} = (x_{\mathrm{S}}, y_{\mathrm{S}}, z_{\mathrm{S}})\) 为 \(R^{3}\) 中满足 \(z_{N} > \sup(r_{i})_{i \in I}\) 且 \(z_{\mathrm{S}} < -\sup(r_{i})_{i \in I}\) 的点。

\begin{enumerate}
\def\labelenumi{\alph{enumi})}
\tightlist
\item
  称道路 \(c\) 在 \(\mathbf{R}^3-K\) 中且起点为 N、终点为 S 为直的，若存在相邻的道路 \(c_1\) 和 \(c_2\) 位于 \(\mathbf{R}^2\) 中，使得 \(c\) 是以下道路的拼接：\(t \mapsto (c_1(t), z_{\mathrm{N}})\)、\(t \mapsto (c_1(1), (1 - t)z_{\mathrm{N}} + tz_{\mathrm{S}})\) 以及 \(t \mapsto (c_2(t), z_{\mathrm{S}})\)。
\end{enumerate}

证明此时点 \(c_1(1)\) 不属于 P；记 \(A_c\) 为包含该点的 \(\mathbf{R}^2 - \mathbf{P}\) 的道路连通分支。证明两条直道路 \(c\) 和 \(c'\) 严格同伦，当且仅当 \(A_c = A_{c'}\)。

   对于 \(\alpha\) 这个 \(\mathbf{R}^2-\mathbf{P}\) 的道路连通分支，我们选取一条道路 \(c(\alpha)\)，其位于 \(\mathbf{R}^2-\mathbf{K}\) 中，起点为 N、终点为 S，且为直道路，使得 \(\mathrm{A}_{c(\alpha)} = \alpha\)。

\begin{enumerate}
\def\labelenumi{\alph{enumi})}
\setcounter{enumi}{1}
\item
   设 \(i \in I\) 满足 \(m_{i} = 0\)，且 \(\overline{\varphi}_{i}^{-1}(K \cap C_{i})\) 是练习 23 中定义的集合 L；有 \(\overline{\varphi}_{i}^{-1}(P \times \{0\} \cap C_{i}) = Q \times \{0\}\)。对所有 \(j \in Z/4Z\)，令 \(\alpha_{i,j}\) 为 \(R^{2} - P\) 的道路连通分支，它包含 \(\varphi_{i}\) 所映的连通分支 \(A_{j}\)，该分支属于 \(B_{2} - P\)。证明关系式 \([c(\alpha_{i,1})][c(\alpha_{i,4})]^{-1} = [c(\alpha_{i,2})][c(\alpha_{i,3})]^{-1}\) 在 \(\pi_{1}(R^{3} - K, N)\) 中成立。
\item
   令 \(\omega\) 是 \(\pi_0(\mathbf{R}^2-\mathbf{P})\) 中的元素。令 \(\lambda_\omega\colon \mathrm{F}(\pi_0(\mathbf{R}^2-\mathbf{P}))\to \pi_1(\mathbf{R}^3-\mathbf{K},\mathbf{N})\) 为唯一的群同态，对于每个 \(\alpha \in \pi_0(\mathbf{R}^2-\mathbf{P})\)，它把元素 \(x_\alpha\) 映为 \([c(\alpha)][c(\omega)]^{-1}\)。
\end{enumerate}

证明同态 \(\lambda_{\omega}\) 是满射，其核是 \(\mathrm{F}(\pi_0(\mathbf{R}^2 - \mathbf{P}))\) 的包含元素 \(x_{\omega}\) 以及 \(x_{\alpha_{i,1}}^{-1}x_{\alpha_{i,2}}x_{\alpha_{i,3}}^{-1}x_{\alpha_{i,4}}\)（其中 \(i \in \mathrm{I}\) 满足 \(m_i = 0\)）的最小正规子群。

\begin{enumerate}
\def\labelenumi{\alph{enumi})}
\setcounter{enumi}{3}
\item
   对于每条线段 \(s\)，它属于 \(\pi_{0}(\mathrm{K}^{\prime})\)，选取元素 \(\alpha_{s}\) 和 \(\alpha_{s}^{\prime}\)，它们属于 \(\pi_{0}(\mathbf{R}^{2}-\mathbf{P})\)，使得 \(R^{2}-P\) 中闭包包含 \(s\) 的连通分支集合等于 \(\{\alpha_{s},\alpha_{s}^{\prime}\}\)。令 \(\mu\colon\mathrm{F}(\pi_{0}(\mathrm{K}^{\prime}))\to\pi_{1}(\mathbf{R}^{3}-\mathbf{K},\mathbf{N})\) 为唯一的群同态，对于 \(s\in\pi_{0}(\mathrm{K}^{\prime})\)，它把元素 \(x_{s}\) 映为 \([c(\alpha_{s})][c(\alpha_{s}^{\prime})]^{-1}\)。证明 \(\mu\) 是满射。对于 \(i\in I\)，其中 \(m_{i}=0\)，以及 \(j\in\{1,2,3,4\}\)，令 \(s_{i,j}\) 为 \(\pi_{0}(\mathrm{K}^{\prime})\) 中满足 \(\varphi_{i}(a_{j})\) 与 \(s_{i,j}\) 相交的唯一线段。证明存在元素 \(u_{i},v_{i},w_{i}\)，它们属于 \(\{-1,1\}\)，使得 \(\mu(x_{s_{i,4}})=\mu(x_{s_{i,2}}^{w_{i}})\) 且 \(\mu(x_{s_{i,3}})=\mu(x_{s_{i,2}}^{-u_{i}}x_{s_{i,1}}^{v_{i}}x_{s_{i,2}}^{u_{i}})\)。证明 \(\mu\) 的核是 \(\mathrm{F}(\pi_{0}(\mathrm{K}^{\prime}))\) 的包含元素 \(x_{s_{i,3}}x_{s_{i,2}}^{-u_{i}}x_{s_{i,1}}^{-v_{i}}x_{s_{i,2}}^{u_{i}}\) 和 \(x_{s_{i,4}}x_{s_{i,2}}^{-w_{i}}\)（对 \(i\in I\)）的最小正规子群。
\item
  Set
\end{enumerate}

\[
k = \operatorname{Card} (\pi_ {0} (\mathrm{K})) - \operatorname{Card} (\{i \in \mathrm{I} \mid m _ {i} > 0 \}) + \frac {1}{2} \sum_ {i \in \mathrm{I}} m _ {i}.
\]

   证明 \(k \in N\)。证明 \(\pi_{1}(\mathbf{R}^{3} - \mathbf{K})\) 对其导出群的商同构于 \(Z^{k}\)。

\begin{enumerate}
\def\labelenumi{\arabic{enumi})}
\setcounter{enumi}{25}
\tightlist
\item
  在数值平面 \(\mathbf{R}^2\) 中，令 \(\mathrm{D}_0\) 为中心为 \((0,0)\)、半径为 1 的圆盘，令 \(\mathrm{C}_0\) 为其边界；置 \(a_0 = (-1,0)\)。令 \(g\) 是满足 \(\geqslant 1\) 的自然整数，令 \(r\) 是严格正实数，令 \((v_1,\ldots,v_g)\) 是实数序列，满足 \(r < \inf (v_{1} + 1,(v_{2} - v_{1}) / 2,\dots,(v_{g} - v_{g - 1}) / 2,1 - v_{g})\)。对于所有 \(j\in \{1,\ldots,g\}\)，记 \(\mathrm{D}_j\) 为中心为 \((0,v_j)\) 的闭圆盘
\end{enumerate}

半径为 \(r\)，并令 \(\mathrm{C}_j\) 为中心为 \((0,v_j)\)、半径为 \(r\) 的圆；再置 \(a_{j}=(-r,v_{j})\)。令 X 为空间 \(\mathrm{D}=\bigcup_{j=1}^{g}\mathring{\mathrm{D}}_{j}\)。

  对每个 \(j \in \{1, \ldots, j\}\)，令 \(u_j\) 为 \(\varpi(\mathrm{X})\) 中一条起点为 \(a_0\)、终点为 \(a_j\) 且像为线段 \([a_0, a_j]\) 的道路的道路类；也记 \(c_j \in \pi_1(\mathrm{X}, a_j)\) 为由 \(t \mapsto (-r \cos(2\pi t), -r \sin(2\pi t) + v_j)\) 给出的圈的道路类，并令 \(\gamma_j = u_j c_j u_j^{-1}\)。

\begin{enumerate}
\def\labelenumi{\alph{enumi})}
\item
  令 \(\gamma_{0}\) 为 X 中由 \(t \mapsto (-\cos(2\pi t), -\sin(2\pi t))\) 给出的圈的道路类。证明有 \(\gamma_{0} = \gamma_{1} \ldots \gamma_{g}\)。
\item
  证明从以 \(x_{1},\ldots,x_{j}\) 为生成元的自由群到 \(\pi_{1}(X,a_{0})\) 的唯一同态（将 \(x_{j}\) 映为 \(\gamma_{j}\)）是群同构。（将范坎彭定理应用于 X 与半平面 \(R_{+}\times R\) 和 \(R_{-}\times R\) 的交所定义的覆盖。）
\item
  令 \(f \colon \mathrm{X} \to \mathbf{R}_+\) 是连续函数，满足 \(f^{-1}(0) = \bigcup_{j=0}^{g} \mathrm{C}_j\)，令 Y 为 \(\mathbf{R}^3\) 中的点 \((x, y, z)\) 所成的子空间，这些点满足 \(z^2 = f(x, y)\)。记 \(i_+\)（分别为 \(i_-\)）为从 X 到 Y 的连续映射，分别由 \((x, y) \mapsto (x, y, \sqrt{f(x, y)})\)（分别由 \((x, y) \mapsto (x, y, -\sqrt{f(x, y)})\)）给出。对于 X 中的每个道路类 \(\gamma\)，在 \(\varpi(\mathrm{Y})\) 中定义道路类 \(\gamma^+ = (i_+)_{*}(\gamma)\) 和 \(\gamma^- = (i_-)_*(\gamma)\)。
\end{enumerate}

对于每个 \(j \in \{1, \ldots, g\}\)，最后在 \(\pi_1(Y, a_0)\) 中定义元素 \(\delta_j = u_j^+ u_j^-\)。证明有关系 \(\gamma_j^+ \delta_j = \delta_j \gamma_j^-\)（\(1 \leqslant j \leqslant g\)），以及 \(\gamma_1^+ \ldots \gamma_g^+ = \gamma_1^- \ldots \gamma_g^-\)。

\begin{enumerate}
\def\labelenumi{\alph{enumi})}
\setcounter{enumi}{3}
\tightlist
\item
  令 \(\varphi\) 为从具有 2g 个生成元 \(x_{1},\ldots ,x_{g},y_{1},\ldots ,y_{g}\) 的自由群到 \(\pi_1(\mathrm{Y},a_0)\) 的唯一同态，它把 \(x_{j}\) 映为 \(\gamma_j^+\)，把 \(y_{j}\) 映为 \(\delta_{j}\)。证明同态 \(\varphi\) 是满射，且其核是包含元素
\end{enumerate}

\[
(x _ {1} ^ {- 1} y _ {1} x _ {1}) (x _ {2} ^ {- 1} y _ {2} x _ {2}) \dots (x _ {g} ^ {- 1} y _ {g} x _ {g}) (y _ {1} \dots y _ {g}) ^ {- 1}.
\]

（将范坎彭定理应用于 Y 与半空间 \(R^{2} \times R_{+}\) 和 \(R^{2} \times R_{-}\) 的交所定义的覆盖。）

\begin{enumerate}
\def\labelenumi{\alph{enumi})}
\setcounter{enumi}{4}
\tightlist
\item
  令 \(\varphi'\) 为从具有 2g 个生成元 \(x_1, \ldots, x_g, y_1, \ldots, y_g\) 的自由群到 \(\pi_1(Y, a_0)\) 的唯一同态，对于每个 \(j\)，它把 \(x_j\) 映为
\end{enumerate}

\[
(\gamma_ {1} ^ {+} \dots \gamma_ {j - 1} ^ {+}) ^ {- 1} \delta_ {j} (\gamma_ {1} ^ {+} \dots \gamma_ {j - 1} ^ {+})
\]

并把 \(y_{j}\) 映为 \(\gamma_{j}^{+}\)。证明同态 \(\varphi'\) 是满射，且其核是包含元素

\[
\left(x _ {1} ^ {- 1} y _ {1} ^ {- 1} x _ {1} y _ {1}\right) \dots \left(x _ {g} ^ {- 1} y _ {g} ^ {- 1} x _ {g} y _ {g}\right).
\]

\begin{enumerate}
\def\labelenumi{\arabic{enumi})}
\setcounter{enumi}{26}
\tightlist
\item
  令 D 为 C 中的单位圆盘，令 X 为 D 的一个子空间，且是 \(S_{1}\) 的一个邻域。令 R 为 X 上的最粗等价关系：对于 \(u, v \in S_{1}\)，当且仅当点 \(f(u)\) 和点 \(f(v)\) 满足 \(f(u) = f(v)\) 时，识别它们。令 \(p\colon X \to X / R\) 为典范满射。令 \(T\) 为紧拓扑空间，令 \(f\colon S_1 \to T\) 为连续满射。置 \(a = f(1)\)，并记 \(\alpha \in \pi_1(T, a)\) 为由 \(t \mapsto f(e^{2\pi it})\) 给出的圈的道路类，它位于 \(T\) 中。
\end{enumerate}

\begin{enumerate}
\def\labelenumi{\alph{enumi})}
\item
  证明，存在唯一的从 T 到 X/R 的映射 \(\sigma\)，使得 \(\sigma \circ f = p|_{\mathbf{S}_1}\)。证明 \(\sigma\) 将 T 同胚到其像 \(p(\mathbf{S}_1)\)。
\item
  令 r 是满足 0 \textless{} r \textless{} 1 的实数；假设 X 是由 \(z \in D\) 且满足 \(r \leqslant |z| \leqslant 1\) 的点所成的集合。证明 X/R 同胚于 f 的映射柱。推出同态 \(\sigma_{*}\) 是从 \(\pi_{1}(T, a)\) 到 \(\pi_{1}(X/R, p(1))\) 的群同构。
\item
  假设 X = D。证明 X/R 同胚于 f 的映射锥。推出群同态 \(\sigma_{*}\colon\pi_{1}(\mathrm{T},a)\to\pi_{1}(\mathrm{X}/\mathrm{R},p(1))\) 是满射，其核是 \(\pi_{1}(\mathrm{T},a)\) 中包含 \(\alpha\) 的最小正规子群。
\item
  令 r 是满足 0 \textless{} r \textless{} 1 的实数；令 \(X_{1}\) 为由 \(z \in X\) 且满足 \(|z| \leqslant r\) 的 z 所成的集合，令 \(X_{2}\) 为由 \(z \in D\) 且满足 \(r \leqslant |z| \leqslant 1\) 的 z 所成的集合；假设 \(X_{2} \subset X\)。令 \(\beta \in \pi_{1}(X_{2}, r)\) 为由 \(t \mapsto re^{2\pi it}\) 给出的圈的道路类；令 \(\delta \in \varpi_{p(r), p(1)}\) 为由 \(t \mapsto p((1 - t)r + t)\) 给出的道路的道路类。
\end{enumerate}

  证明，存在唯一的群同态 \(\mu\)，从 \(\pi_1(\mathrm{T},a)*\pi_1(\mathrm{X}_2,r)\) 到 \(\pi_1(\mathrm{X} / \mathrm{R},p(1))\)，它与同态 \(\sigma_*\) 在 \(\pi_1(\mathrm{T},a)\) 上一致，并与同态 \(v\mapsto \delta^{-1}p_{*}(v)\delta\) 在 \(\pi_1(\mathrm{X}_2,p(r))\) 上一致。证明 \(\mu\) 是满射，且其核是包含 \(\alpha \beta^{-1}\) 的最小正规子群。

\begin{enumerate}
\def\labelenumi{\arabic{enumi})}
\setcounter{enumi}{27}
\tightlist
\item
  设同伦 \(\sigma: \mathrm{X} \times \mathbf{I} \to \mathrm{X}\) 为同痕，若对每个 \(t \in \mathbf{I}\)，映射 \(x \mapsto \sigma(x, t)\) 是由 \(\sigma\) 诱导的从 X 到自身的同胚。
\end{enumerate}

\begin{enumerate}
\def\labelenumi{\alph{enumi})}
\item
  令 m 是自然整数，令 \(f: X \to R^{n}\) 为连续映射。证明，存在一个同痕，其起点是 \(X \times R^{n}\) 的恒等映射，终点是从 \(X \times R^{n}\) 到 \(X \times R^{n}\) 的映射 \((x, y) \mapsto (x, y - f(x))\)。
\item
  令 A、B 是 X 的子空间；若存在同痕 \(\sigma\colon\mathbf{X}\times\mathbf{I}\to\mathbf{X}\)，其起点是 X 上的恒等映射，且满足 \(f(\mathbf{A}\times\{1\})=\mathbf{B}\)，则称 A 与 B 同痕。
\end{enumerate}

证明关系“ A 与 B 同痕”是在 X 的子空间集合（分别在 X 的闭子空间集合）上的等价关系。

\begin{enumerate}
\def\labelenumi{\alph{enumi})}
\setcounter{enumi}{2}
\tightlist
\item
  设 m、n 是自然数，A 是 \(R^{m}\) 的闭子空间，B 是 \(R^{n}\) 的闭子空间。假设 A 与 B 同胚。证明 \(A \times \{0\}\) 和 \(\{0\} \times B\) 是 \(R^{m+n}\) 的同痕子空间。
\end{enumerate}

\begin{enumerate}
\def\labelenumi{\arabic{enumi})}
\setcounter{enumi}{28}
\tightlist
\item
  设 A 是 \(R^{2}\) 的闭子集，并且同胚于 R 的一个闭子集 B。
\end{enumerate}

\begin{enumerate}
\def\labelenumi{\alph{enumi})}
\item
  证明 \(A \times \{0\}\) 和 \(\{(0,0)\} \times B\) 是 \(R^{3}\) 的同痕子空间。
\item
  假设 \(B \neq R\)；证明 \(R^{2}-A\) 道路连通。
\item
  假设 B = R；证明 \(R^{2}-A\) 恰有两个道路连通分支，且 A 是每个分支的边界。
\item
  设 A 是 \(S_{2}\) 中同胚于 \(S_{1}\) 的子集；证明 \(S_{2} - A\) 恰有两个道路连通分支，且 A 是每个分支的边界。
\item
  设 A 是 \(R^{2}\) 中同胚于 \(S_{1}\) 的子集；证明 \(R^{2}-A\) 恰有两个道路连通分支，且 A 是每个分支的边界。
\end{enumerate}

\begin{enumerate}
\def\labelenumi{\arabic{enumi})}
\setcounter{enumi}{29}
\tightlist
\item
  令 \(v = (v_{1},\ldots ,v_{n})\colon \mathbf{R}^{n}\to \mathbf{R}^{n}\) 是 \(\mathrm{C}^1\) 类映射，令 \(\delta\) 是满足 \(>0\) 的实数；作如下假设：
\end{enumerate}

- \(v\) 的支集包含于 \([- \delta, 1 + \delta] \times \mathring{\mathbf{B}}_{n-1}\) 中；

- 对每个 \(x \in [0, 1 - \delta[\)，都有 \(v_1(x, 0, \ldots, 0) > 0\)，且 \(v_1(1, 0, \ldots, 0) = 0\)；

- 对每个 \(x \in [0,1]\)，都有 \(v_{2}(x,0,\ldots,0) = \cdots = v_{n}(x,0,\ldots,0) = 0\)。

  令 \(\Phi\colon\mathbf{R}^{n}\times\mathbf{R}\to\mathbf{R}^{n}\) 为映射 \((x,t)\mapsto v(x)\) 的积分流（VAR，9.1.3）；对 \(t\in\mathbf{R}\)，令 \(\Phi_{t}\) 为映射 \(x\mapsto\Phi(x,t)\)。

\begin{enumerate}
\def\labelenumi{\alph{enumi})}
\item
  证明存在实数 \(\tau\)，使得 \(\Phi_{\tau}([0,1] \times \{0\}) \subset [1 - \delta, 1] \times \{0\}\)。
\item
  令 \(j\colon[-\delta,1+\delta]\times\mathbf{B}_{n-1}\to\mathbf{R}^{n}\) 是连续单射映射。对每个 \(t\in[0,1]\)，置 \(\mathrm{A}_{t}=j([t,1]\times\{0\})\)。证明对每个 \(t\in]0,1[\)，子空间 \(A_{0}\) 和 \(A_{t}\) 同痕。
\item
  设 P 是 \(\mathbf{R}^{n}\) 中同胚于 \([0,1]\) 且为有限族线段之并的子集。证明 P 与一条线段同痕。
\end{enumerate}

\begin{enumerate}
\def\labelenumi{\arabic{enumi})}
\setcounter{enumi}{30}
\tightlist
\item
  设 \(f = (f_1, \ldots, f_n) \colon \mathbf{I} \to \mathbf{R}^n\) 是 \(\mathbf{C}^1\) 类映射，满足 \(f'(t) \neq 0\) 对所有 \(t \in \mathbf{I}\) 都成立。
\end{enumerate}

\begin{enumerate}
\def\labelenumi{\alph{enumi})}
\item
  假设 \(f_{1}^{\prime}(t) > 0\) 对所有 \(t \in I\) 都成立。对每个 \(j \in \{2, \ldots, n\}\)，证明存在连续映射 \(g_{j} \colon R \to R\)，使得 \(f_{j}(x) = g_{j}(f_{1}(x))\) 对所有 \(x \in I\) 都成立。推出 \(f(\mathbf{I})\) 与一条线段同痕。
\item
  证明存在实数 \(\delta > 0\) 和连续单射映射 \(j\colon [-\delta, 1 + \delta] \times \mathbf{B}_{n-1} \to \mathbf{R}^n\)，使得 \(j(t, 0) = f(t)\) 对所有 \(t \in \mathbf{I} .\) 都成立。由练习 30 推出，对每个 \(\tau \in [0, 1[\)，集合 \(f(\mathbf{I})\) 与 \(f([\tau, 1])\) 同痕。
\item
  证明 \(f(\mathbf{I})\) 与一条线段同痕。
\end{enumerate}

\begin{enumerate}
\def\labelenumi{\arabic{enumi})}
\setcounter{enumi}{31}
\tightlist
\item
  设 A 是 \(\mathbf{R}^3\) 中同胚于 \([0,1]\) 的闭子集。A 的端点是 A 中的两个点 \(p\)，满足 \(\mathrm{A} - \{p\}\) 连通（参见 III，第 280 页，定理 2）；A 的其他点称为内部点。
\end{enumerate}

若 \(p\) 是 A 的一点；称 A 在 \(p\) 处是驯的，若存在 \(p\) 在 \(\mathbf{R}^3\) 中的邻域 V，以及从 V 到 \(\mathbf{R}^3\) 的单位球的同胚，它把 \(\mathrm{A} \cap \mathrm{V}\) 映为一条线段。

\begin{enumerate}
\def\labelenumi{\alph{enumi})}
\item
  设 \((\mathrm{V}_{n})\) 是 \(R^{3}\) 中 p 的邻域的递减序列，满足 \(\bigcap_{n}\mathrm{V}_{n}=\{p\}\)；对每个 n，令 \(a_{n}\) 是 \(V_{n}-A\) 中的一点，并记 \(j_{n}\colon\pi_{1}(\mathrm{V}_{n}-\mathrm{A},a)\to\pi_{1}(\mathrm{V}_{1}-\mathrm{A},a)\) 为由 \(V_{n}\) 包含于 \(V_{1}\) 所诱导的群同态。若 p 是 A 的端点（分别为内部点），证明存在整数 m，使得对每个 \(n\geqslant m\)，同态 \(j_{n}\) 的像归结为单位元（分别为阿贝尔群）。
\item
  假设 A 在每一点处都是驯的。证明存在从 \(R^{3}\) 到自身的同胚，将 A 映为一条线段。推出 \(R^{3}-A\) 连通且道路单连通。
\end{enumerate}

\begin{enumerate}
\def\labelenumi{\arabic{enumi})}
\setcounter{enumi}{32}
\tightlist
\item
  \begin{enumerate}
  \def\labelenumii{(\arabic{enumii})}
  \setcounter{enumii}{5}
  \tightlist
  \item
     设 L 是 \(\mathbf{R}^3\) 的子空间，定义见练习 23，其中取 \(u = (1 - i)/\sqrt{2}\) 和 \(v = (1 + i)/\sqrt{2}\)。令 \(a\)、\(r\) 是满足 \(0 < 2r < a < 1 - 2r\) 的实数。置 \(p_1 = (0, -a)\)、\(p_2 = (0, 0)\)、\(p_3 = (0, a)\)；\(q_1 = (-1, -a)\)、\(q_2 = (-1, 0)\)、\(q_3 = (-1, a)\)；\(q_1' = (1, -a)\)、\(q_2' = (1, 0)\)、\(q_3' = (1, a)\)。令 K 为集合 \(p_1 + rL\)、\(p_2 + rL\)、\(p_3 + rL\) 与线段 \([q_1, p_1 - rv]\)、\([q_2, p_2 - ru]\)、\([q_3, p_3 - ru]\)、\([q_1', p_1 + ru]\)、\([q_2', p_3 + ru]\)、\([q_3', p_3 + rv]\) 的并。
  \end{enumerate}
\end{enumerate}

令 \((x_{n})_{n\in \mathbf{Z}}\) 是严格递增的实数族，满足 \(\lim_{n\to -\infty}x_n = -1\) 且 \(\lim_{n\to +\infty}x_n = 1\)；对 \(n\in \mathbf{Z}\)，置 \(z_{n} = y_{n} = 1 - |x_{n}|\)。对 \(n\in \mathbf{Z}\)，令 \(f_{n}\) 为从 \([-1,1]^{3}\) 到 \(\mathbf{R}^3\) 的映射，由下式给出：

\[
f _ {n} (x, y, z) = \frac {1 - x}{2} (x _ {n}, y y _ {n}, z z _ {n}) + \frac {1 + x}{2} (x _ {n + 1}, y y _ {n + 1}, z z _ {n + 1})
\]

，并置 \(A_{n}=f_{n}(K)\) 和 \(B_{n}=f_{n}([-1,1]^{3})\)。令 A 为族 \((\mathrm{A}_{n})_{n\in\mathbf{Z}}\) 的并。

\begin{enumerate}
\def\labelenumi{\alph{enumi})}
\item
  证明 A 同胚于 [0, 1]。

(6) 此例出自 E. ARTIN 与 R. FOX 的文章《三维空间中的一些野胞与球面》，载于《数学年刊》49 (1948)，第 979–990 页。

\item
对于每个自然整数 \(m\)，置
\end{enumerate}

\[
\mathrm{A} _ {m} ^ {\prime} = \bigcup_ {n <   - m} \mathrm{B} _ {n} \cup \bigcup_ {- m \leqslant n \leqslant m} \mathrm{A} _ {m} \cup \bigcup_ {n > m} \mathrm{B} _ {n}.
\]

证明 \(R^{3}-A_{m}^{\prime}\) 道路连通，且其基本群由元素 \(a_{n}, b_{n}, c_{n}\)（其中 \(n \in Z\) 且 \(-m \leqslant n \leqslant m\)）生成，并满足关系 \(b_{-m} = c_{-m}a_{-m}\)、\(c_{n+1}a_{n+1} = c_{n}c_{n+1}\)、\(c_{n+1}b_{n} = a_{n}c_{n+1}\)、\(b_{n}c_{n+1} = b_{n+1}b_{n}\)（其中 \(n \in Z\) 且 \(-m \leqslant n < m\)）以及 \(c_{m}a_{m} = b_{m}\)。

\begin{enumerate}
\def\labelenumi{\alph{enumi})}
\setcounter{enumi}{2}
\tightlist
\item
   证明 \(R^{3}-A\) 道路连通，且其基本群由元素 \(a_{n}, b_{n}, c_{n}\)（\(n \in Z\)）生成，并满足关系
\end{enumerate}

\[
b _ {n} = c _ {n} a _ {n}, \quad c _ {n + 1} a _ {n + 1} = c _ {n} c _ {n + 1}, \quad c _ {n + 1} b _ {n} = a _ {n} c _ {n + 1}, \quad b _ {n} c _ {n + 1} = b _ {n + 1} b _ {n}
\]

 当 \(n \in Z\) 时。

d) 证明存在从 \(\pi_1(\mathbf{R}^3 -\mathbf{A})\) 到群 \(\mathfrak{S}_5\) 的同态，它把 \(c_{n}\) 的道路类映为置换 \((1,2,3,4,5)\mapsto (2,3,4,5,1)\)（当 \(n\) 为奇数），并映为置换 \((1,2,3,4,5)\mapsto (4,3,5,2,1)\)（当 \(n\) 为偶数）。推出 \(\mathbf{R}^3 -\mathbf{A}\) 不是道路单连通的。

\begin{enumerate}
\def\labelenumi{\arabic{enumi})}
\setcounter{enumi}{33}
\tightlist
\item
  设 X 是道路连通的拓扑空间，G 是连续左作用于 X 的离散群。令 M 是 X 的开子集，满足 G·M = X。按照 I，第 136 页命题 10 定义集合 S、T、群 F 及群同态 \(\varphi\colon\mathrm{F}\to\mathrm{G}\)。令 \(a\) 是 M 中的一点。
\end{enumerate}

\begin{enumerate}
\def\labelenumi{\alph{enumi})}
\item
  设 \(c\colon\mathbf{I}\to\mathbf{X}\) 是以 \(a\) 为基点的圈。证明存在 F 中唯一的元素 \(g(c)\)，满足以下性质：对每个整数 \(n\) 以及 S 中元素的每个序列 \((s_{1},\ldots,s_{n})\)，若 \(c([(k-1)/n,k/n])\subset s_{k}\mathrm{M}\) 对每个整数 \(k\in\{1,\ldots,n\}\) 都成立，则 \(g(c)=x_{s_{n}^{-1}s_{1}}x_{s_{1}^{-1}s_{2}}\ldots x_{s_{n-1}^{-1}s_{n}}\)。
\item
  证明存在唯一的群同态 \(\gamma\)，从 \(\pi_1(\mathrm{X},a)\) 到 F，使得 \(\gamma ([c]) = g(c)\)，其中 \(c\) 是 X 中以 \(a\) 为基点的每个圈。
\item
  证明 \(\varphi\) 的核等于同态 \(\gamma\) 的像。
\item
  证明 \(\gamma\) 的核由各群 \(\pi_{1}(\mathrm{M}\cup s\cdot\mathrm{M},a)\)（\(s\in S\)）的像的并生成。
\end{enumerate}

